\chapter{MuStar: Semantic-to-Kinetic Compilation}

In the \texttt{uni*} ecosystem, there is a fundamental gap between semantic process intent generated by intelligence layers and the physical motion required by the execution substrate. \textbf{MuStar} bridges this gap. It compiles intent into geometric motion.

\section{POWL Intent}

When DTEAM dictates an operational change, it uses the Partially Ordered Workflow Language (POWL) to express partial-order intent. MuStar takes this high-level business logic and begins the compilation process.

\section{Ontology Constraints}

MuStar binds the POWL intent against the predefined ontology constraints of the ecosystem, ensuring that the semantic request is structurally valid within the target domain.

\section{Active Scopes over $64^2$}

The compiler evaluates the geometric boundaries within the attention surface ($64^2$) where motion is permitted, generating the active geometric scopes for the operation.

\section{Motion Packets over $64^3$}

MuStar translates the operation into specific bit-level target states within the $64^3$ operational globe, generating \textbf{Candidate Masks}. These represent the exact geometric coordinates and bit-level masks required by the physical substrate.

\section{Work-Tier Hints over $8^n$}

MuStar generates pre-calculated geometric bounds for the $8^n$ execution loops, providing work-tier hints to the downstream execution engine.

\section{Proof Obligations}

Finally, MuStar attaches cryptographic or structural assertions to the motion packet. These are the proof obligations that \texttt{unios} must successfully verify before granting admission to the physical substrate. MuStar itself does not execute, admit, or prove authority; it strictly compiles the semantic intent into this \texttt{unios}-admissible packet.